\newpage
\section{模型的建立与求解}

\subsection{数据预处理}
\begin{enumerate}[label=(\arabic*), itemindent=0pt, labelindent=\parindent, labelwidth=2em, labelsep=5pt, leftmargin=*]
\item[1] 缺失值检测：采用3\sigma准则识别异常值，线性插值补全
\item[2] 归一化：Min-Max标准化 x*=(x-x_min)/(x_max-x_min)
\item[3] 特征工程：构造小时特征、星期特征、节假日标志
\end{enumerate}

\subsection{问题一模型：负荷预测}
\subsubsection{STL分解}
STL将时间序列分解为：\F_t = T_t + S_t + R_t
其中T_t为趋势项，S_t为季节项（周期tau=24小时），R_t为残差项。

\subsubsection{Transformer预测}
Self-Attention机制：\\text{Attention}(Q,K,V) = \text{softmax}(\frac{QK^T}{\sqrt{d_k}})V
多头注意力并行计算h个子空间的注意力。位置编码采用正弦函数。

\subsubsection{预测结果}
\begin{table}[H]
\centering
\caption{不同模型预测性能对比}
\begin{tabular}{c c c c}
\toprule
模型 & MAE & RMSE & MAPE(%) 
\midrule
ARIMA & 0.052 & 0.067 & 6.1 
LSTM & 0.041 & 0.053 & 4.8 
\textbf{STL-Transformer} & \textbf{0.028} & \textbf{0.036} & \textbf{3.5} 
\bottomrule
\end{tabular}
\end{table}

\subsection{问题二模型：多能互补调度}
\subsubsection{目标函数}
\begin{align}
\min Z_1 &= \sum_{t=1}^{T}[c_t^{\text{grid}}P_t^{\text{grid}}+c_{\text{loss}}|P_t^{\text{bat}}|] 
\min Z_2 &= \sum_{t=1}^{T}[\lambda_{\text{carbon}}P_t^{\text{grid}}+\lambda_{\text{pv}}(P_t^{\text{pv,max}}-P_t^{\text{pv}})] 
\end{align}

\subsubsection{约束条件}
\begin{align}
\text{功率平衡：} & P_t^{\text{load}} = P_t^{\text{pv}}+P_t^{\text{grid}}+P_t^{\text{bat}}, \quad \forall t 
\text{储能SOC：} & E_{\min}\leq E_t^{\text{bat}}\leq E_{\max}, \quad \forall t 
\text{爬坡约束：} & |P_t^{\text{bat}}-P_{t-1}^{\text{bat}}|\leq R^{\max}, \quad \forall t 
\end{align}

\subsection{NSGA-II求解}
种群大小N=100，最大迭代G=200，交叉概率pc=0.9，变异概率pm=0.1。得到28个非支配解，成本降低15.3%，碳排放减少22.1%。